% Towards Continual Learning for SmaTables — working draft.
% Build: latexmk -pdf main.tex   (see paper/README.md). Venue TBD; LNCS format assumed.
\documentclass[runningheads]{llncs}

\usepackage[T1]{fontenc}
\usepackage[utf8]{inputenc}
\usepackage{graphicx}
\usepackage{booktabs}
\usepackage{siunitx}
\usepackage{xcolor}
\usepackage{hyperref}
\usepackage{orcidlink}

% Draft helpers — remove before submission.
\newcommand{\todo}[1]{\textcolor{red}{[TODO: #1]}}
\newcommand{\result}[1]{\textcolor{blue}{[RESULT: #1]}}

\begin{document}

\title{Towards Continual Learning for SmaTables: Replay-Buffer Sizing and On-Device Fine-Tuning for Edge-MCU Gesture Recognition}
\titlerunning{Continual Learning for SmaTables on an Edge MCU}

\author{Leo Buron \and Florian Hettstedt \and Andreas Erbslöh \and Gregor Schiele}
\authorrunning{L. Buron et al.}
\institute{Intelligent Embedded Systems, University of Duisburg-Essen, Germany\\
\email{\todo{institutional e-mail addresses}}}

\maketitle

\begin{abstract}
\todo{Write last. Skeleton: SmaTable surface-vibration gesture classifier; all prior work kept inference and training on a PC; we report the first on-RP2350 measurements of both inference and training; on that deployment we quantify (i) how much per-class replay memory prevents forgetting during on-device fine-tuning to a new user (random vs.\ herding exemplars vs.\ PPCA generative replay at matched bytes), (ii) how many new-user samples personalisation needs (LOSO), (iii) how many samples from how many contributors a new gesture needs, (iv) run-to-run stability across seeds, decomposed into init vs.\ data order, (v) whether additive-noise augmentation reduces those budgets. One headline number per RQ.}
\keywords{Continual learning \and On-device training \and Replay \and Microcontroller \and Gesture recognition \and TinyML}
\end{abstract}

\section{Introduction}
\label{sec:intro}

% Outline (from docs/paper-plan.md "Why this paper, why now"):
% 1. SmaTable: hidden piezo sensors in a table surface, privacy-by-design gesture input.
% 2. Deployment exposes three continual-learning questions the literature does not answer
%    for this system: user adaptation budget, new-gesture budget, replay budget.
% 3. Hettstedt et al. 2026 established the DAQ chain on the RP2350 but inference and
%    training stayed on a PC; Yoshida 2023a/b are PC-class and need calibration data to
%    leave the device.
% 4. This paper: first on-RP2350 inference + training measurements, and quantitative
%    sample-budget / stability envelopes for three CL scenarios.

\todo{Motivation paragraph: SmaTable concept and why on-device training matters for it (privacy: calibration data must not leave the table).}

\todo{Gap paragraph: prior SmaTable work \cite{hettstedt2026smatable,yoshida2023a,yoshida2023b} keeps training on a PC; no on-device numbers exist.}

\paragraph{Contributions.}
\begin{enumerate}
  \item \textbf{Replay sizing curve (RQ1).} Accuracy on the original users as a function of per-class replay memory under three replay mechanisms at matched bytes per class --- random exemplars, herding exemplars~\cite{rebuffi2017icarl}, and PPCA generative replay~\cite{tipping1999ppca} --- plus a no-replay floor, measured on-device on the RP2350.
  \item \textbf{User-adaptation sample curve (RQ2).} Accuracy on an unseen user as a function of the number of that user's samples used for on-device fine-tuning, under leave-one-subject-out cross-validation, reporting target-user gain, original-user retention and forward transfer~\cite{lopezpaz2017gem}.
  \item \textbf{New-gesture sample curve (RQ3).} Accuracy on one added gesture class as a function of samples per contributor $\times$ number of contributors, while bounding degradation on the pre-existing classes.
  \item \textbf{On-device training stability (RQ4).} Run-to-run variance of final accuracy across seeds, decomposed into new-parameter initialisation and on-device data order.
  \item \textbf{Augmentation effect (RQ5).} Whether additive Gaussian noise reduces the sample budgets of RQ1/RQ2 at matched accuracy.
\end{enumerate}

\todo{One sentence on what is explicitly \emph{not} claimed: framework additions (Conv1d, LayerNorm, GroupNorm, PPCA replay in the ODT library) are enablers; architecture is fixed; no QPU.}

\section{Related Work}
\label{sec:related}

\paragraph{SmaTable and surface-vibration gesture recognition.}
\todo{Hettstedt et al. 2026 \cite{hettstedt2026smatable}: RP2350 + 9-channel piezo AFE, 4 of 9 channels used, 1\,kHz sampling, 15 users $\times$ 10 sessions $\times$ 10 repetitions $\times$ 6 gestures = 9\,000 events; 1D-SepCNN with 8\,722 parameters, 97.32\,\% 5-fold top-1 offline. Only the DAQ stage runs on the MCU. Yoshida 2023a/b (NAIST) \cite{yoshida2023a,yoshida2023b}: PC-class; 0.90 accuracy under LOPOCV with one calibration session, calibration data leaves the device.}

\paragraph{Replay-based continual learning.}
\todo{Catastrophic forgetting \cite{mccloskey1989catastrophic}; exemplar replay and herding selection \cite{rebuffi2017icarl}; GEM's backward/forward-transfer metrics \cite{lopezpaz2017gem}; regularisation alternatives EWC \cite{kirkpatrick2017ewc} and LwF \cite{li2018lwf} (out of scope here, replay-only). Generative replay with per-class probabilistic PCA \cite{tipping1999ppca}: one Gaussian generator per class, $O((k+1)\,d)$ memory, constant in the number of absorbed domains.}

\paragraph{On-device training on microcontrollers.}
\todo{TinyML on-device training literature to cite (survey + 2--3 systems); position ODT (\texttt{es-ude/OnDeviceTraining}) \cite{odt2026} as a C library whose HOST build is bit-equivalent to the MCU build, which is what makes a supercomputer sweep transferable to the device.}

\section{System and Experimental Setup}
\label{sec:system}

\subsection{Classifier and training framework}
The classifier is the SmaTable \emph{DepthwiseCNN} of Hettstedt et al.~\cite{hettstedt2026smatable}: per block a depthwise \texttt{Conv1d} (kernel $k$, dilation $d$, same padding, no bias), a pointwise $1\times1$ convolution (no bias), \texttt{GroupNorm}(1, $C$), ReLU, \texttt{MaxPool1d}(2) and dropout; the head is adaptive average pooling, a 16-unit linear layer with ReLU and dropout, and a 6-way linear output. Four architecture configurations from the original search are used unchanged (Table~\ref{tab:archs}). \todo{Table~\ref{tab:archs}: trial id, widths, $k$, $d$, dropout, window length $T$, parameter count, PyTorch/Adam LOSO reference accuracy: 2353 (4; 7; 3; 0.1; 1250; 234; 0.618), 3408 (32,8,16,16,8; 23; 1; 0.1; 1250; 3050; 0.904), 3650 (8,8; 5; 2; 0.0; 250; 434; 0.705), 4223 (8,12,8; 7; 3; 0.0; 625; 694; 0.811).}

All training runs use the OnDeviceTraining (ODT) C library~\cite{odt2026} at a pinned revision. The HOST build of the library is bit-equivalent to the MCU build, so hyper-parameter search runs as thousands of HOST processes on a cluster and only the selected configurations are re-executed on the device. Training is FP32 SGD with momentum, cosine learning-rate schedule and best-epoch snapshot; Adam is excluded because its two extra state buffers per parameter do not fit the memory budget. \todo{Cite/describe the parity gates: rebuilt PyTorch model reproduces the reference confusion matrices exactly (V0); prediction and gradient parity vs.\ PyTorch at float32 precision on all four architectures (V1/V2); checkpoint round-trip bitwise (V4); LOSO mean within 3\,pp of the Adam reference (V3).}

\subsection{Hardware}
\todo{RP2350 (dual Cortex-M33, FPU, 520\,KB SRAM) as the primary target; RP2040 (dual Cortex-M0+, no FPU, 264\,KB SRAM) as the portability companion. Timing via the DWT cycle counter; energy via the bench power measurement. Describe the UART dump path used for the HOST$\leftrightarrow$MCU single-step equivalence check (L3).}

\subsection{Dataset and protocol}
\todo{15 subjects, 6 gestures (knock, tap, swipe-down/up/left/right), 8\,999 preprocessed windows of 4 channels; $T$ per configuration; leave-one-subject-out (15 folds). Two-domain protocol for RQ1: domain 0 = pooled original users (fold's train split), domain 1 = held-out user (fold's test split); pretrain on domain-0-fit, snapshot domain-0-held accuracy, fine-tune on domain-1-fit with replay, re-evaluate. BWT $=$ acc$_\text{after}$ $-$ acc$_\text{before}$ (GEM, two domains). Data availability statement: \todo{agree wording with F.~Hettstedt}.}

\subsection{Replay mechanisms and the iso-byte budget}
\todo{Exemplar buffer (random / herding) vs.\ per-class PPCA generator (mean, $k$ basis vectors, $k$ eigenvalues, $\sigma^2$). Budget axis = bytes per class; for PPCA the rank is derived as $k = \text{buffer\_size} - 1$ and the achieved bytes and iso-exemplar count are logged per run rather than assumed.}

\subsection{Search infrastructure}
\todo{Optuna GridSampler with the CV fold as a grid dimension; SQLite-WAL storage on node-local tmpfs; Apptainer container on the amplitUDE Slurm cluster; per-trial timeouts; trial counts per study.}

\section{Results}
\label{sec:results}

\subsection{Memory feasibility on the RP2350}
\label{sec:results-memory}
% AUTOGENERATED from runs/paper-mode-<trial>/ckpt/{memory,manifest}.json (stage-1 host-side probes,
% 3-epoch paper-mode runs, GroupNorm mode, pin 7d7f1d5). Regenerate rather than edit.
\begin{table}[t]
\centering
\caption{Memory feasibility of the four DepthwiseCNN configurations on the RP2350 (520\,KB SRAM). Static budget = parameters + gradients + optimiser state + activations + gradient buffer + I/O + masks, measured by the host-side probes of the C trainer; stack peak measured by painted-stack scan over a full training run. Reference accuracy is the PyTorch/Adam LOSO mean of the original search. \todo{confirm on device}}
\label{tab:memory-feasibility}
\begin{tabular}{lrrrrrrrc}
\toprule
Config & Widths & $k$ & $d$ & $T$ & Params & Ref.\ acc. & Static + stack (KiB) & Fits \\
\midrule
trial-3408 & 32,8,16,16,8 & 23 & 1 & 1250 & 3050 & 0.904 & 1953 & no \\
trial-4223 & 8,12,8 & 7 & 3 & 625 & 694 & 0.811 & 303 & yes \\
trial-3650 & 8,8 & 5 & 2 & 250 & 434 & 0.705 & 109 & yes \\
trial-2353 & 4 & 7 & 3 & 1250 & 234 & 0.618 & 242 & yes \\
\bottomrule
\end{tabular}
\end{table}


Three of the four configurations fit the RP2350's 520\,KB SRAM with headroom; the most accurate configuration (trial-3408, 0.904 reference accuracy) does not: its activation and stack footprint scales with width\,$\times$\,$T$ and exceeds the budget by roughly $3.8\times$. \todo{Confirm the host-side estimates on the device (peak SRAM via painted stack + linker map) before this becomes a claim; keep the host numbers as the estimate column.}

\subsection{Stage 1: LOSO pretraining in C}
\result{Per-architecture LOSO mean $\pm$ std over 15 folds vs.\ the PyTorch/Adam reference; V3 gate. Local pilot (trial-3650, single LR point 0.1, 250 epochs): 0.7205 $\pm$ 0.0898 vs.\ 0.7045 reference.}

\subsection{RQ1: replay sizing at matched memory}
\result{Accuracy on original users vs.\ bytes per class, four arms, error bars from the seed axis; BWT; new-user accuracy as a sanity check. Secondary: PPCA rank sensitivity; absorption chunk-size sensitivity (cost curve).}

\subsection{RQ2: unseen-user sample efficiency}
\result{Target-user accuracy vs.\ $K$; original-user retention; FWT at $K=0$.}

\subsection{RQ3: new-gesture sample efficiency}
\result{New-class accuracy vs.\ samples per contributor $\times$ contributors; old-class retention floor.}

\subsection{RQ4: on-device training stability}
\result{Std-dev across seeds; variance shares init vs.\ data order.}

\subsection{RQ5: noise augmentation}
\result{$\Delta$ accuracy at matched budget, aug on vs.\ off.}

\subsection{On-device measurements}
\result{Inference latency, fine-tune step latency, energy per inference / per epoch, peak SRAM on the RP2350; RP2040 companion numbers if available; HOST$\leftrightarrow$MCU single-step parity (L3).}

\section{Discussion}
\label{sec:discussion}

\todo{What the budgets mean for the SmaTable product: how much replay memory to reserve per class, how many calibration gestures to ask a new user for, how many contributors a new gesture needs.}

\todo{Deployability vs.\ accuracy: the most accurate configuration is undeployable on the RP2350; the 0.811 configuration fits with $\sim$40\,\% headroom.}

\paragraph{Limitations.}
\todo{RQ1 measures exactly one domain transition (original users $\rightarrow$ one new user), so PPCA's constant-memory-across-domains property is not exercised. 15 subjects. Fixed architecture. FP32 only. Energy scope. Replay-only (no EWC/LwF/NCM).}

\section{Forthcoming Work}
\label{sec:forthcoming}

\todo{Joint HW-NAS + CL search with an analytical memory/MAC cost model; QPU variant of the same matrix once QPU integration lands in ODT; environmental-noise adaptation (standalone paper); richer replay selection (gradient-, uncertainty-, loss-based); additional CL methods (EWC-diag, LwF, NCM head); a sequential multi-user domain chain to exercise PPCA's memory advantage.}


\bibliographystyle{splncs04}
\bibliography{references}

\end{document}
